\section[Curvature estimates and the Abelian clamp]{Uniform curvature estimates and the Abelian clamp}
\label{sec:estimates}

We now turn the metric periods and tube estimates of
Section~\ref{sec:geometry} into simultaneous bounds on the parallel line
classes that enter the index calculation of Section~\ref{sec:indices}.

The metric choices in Section~\ref{sec:geometry} will be made for a
\emph{fixed finite problem}.  This means that the number of intervals, the
bits, the bundles, all exceptional evaluations, the charge bound, and the
finite set of reference fillings have been fixed.  Denote this collection
by $\mathcal P$.  Constants with a subscript $\mathcal P$ are allowed to
depend on all of it.  The constants attached to an old cap $W$ below have
a different, explicitly stated, dependence.

We use real two-forms $\beta,D$, with
\[
 b=i\beta,\qquad [\beta/(2\pi)]=\Lambda,
 \qquad [D/(2\pi)]=v.
\]
Thus $iD$ is the difference of the two line curvatures at a reduction,
and the active spinor line has determinant class $K=\Lambda+v$.
Type $R$ means an Abelian connection with zero spinor, with an ordering
of its summands; type $S$ means a parallel splitting with a nonzero
spinor in the active summand.  These component-type letters do not
denote the capped Seifert surface $A$; the full field-type classification
is Definition~\ref{def:indices:types}.  Our curvature norm is normalized
so that, on a region with its inherited relative charge,
\begin{equation}
 \|F_a^0\|_2^2
   =2\|F_a^{0,+}\|_2^2+8\pi^2\kappa.
 \label{eq:estimates:energy-charge}
\end{equation}
The identity is additive under cutting, whether or not the individual
charges are positive.

For a positive bridge put $X_1=-S_j$, $X_2=S_{j+1}$, as in
Section~\ref{sec:geometry}, and define
\begin{equation}
 \begin{gathered}
 \lambda=\Lambda U=-1-e_j+e_{j+1},\qquad u=vU,\\
 H_I=U+\frac14\sum_{i\in I}X_i,\qquad
 z_I=vH_I,\qquad \lambda_I=\Lambda H_I
 \quad(I\subset\{1,2\}).
 \end{gathered}
 \label{eq:estimates:vertices}
\end{equation}
In particular $H_\varnothing=U$.  Direct substitution gives
\begin{equation}
 \lambda_1=-\frac32+e_{j+1},\qquad
 \lambda_2=-\frac12-e_j,\qquad
 \lambda_{12}=-1,
 \qquad u\equiv\lambda\pmod2.
 \label{eq:estimates:vertex-values}
\end{equation}
A whole side is a side on which the sphere test is still the whole
sphere test.  A separated $N_i$ with its specified Abelian reference
filling counts as whole; a side in the transition to the two-piece
$J_i$ rule need not be used.

\begin{theorem}[The Abelian clamp]
\label{thm:estimates:clamp}
For the finite problem $\mathcal P$ and a fixed upper bound on total
topological charge, the metrics and fixed line connections of
Theorem~\ref{thm:geometry:family} can be chosen with all artificial
isolation lengths finite so that the following assertions hold
simultaneously, including at every genuine face.
\begin{enumerate}[label=\textup{(\roman*)}]
\item On an $R$ or $S$ component containing a positive bridge, either
 $u=0$, or there is a nonempty subset $I$ of its adjacent whole sides
 such that
 \begin{equation}
   \operatorname{sign}(u)z_I\leq0\quad(R),
   \qquad
   \operatorname{sign}(u)z_I<|\lambda_I|\quad(S).
   \label{eq:estimates:clamp}
 \end{equation}
\item No $R$ component contains an outside cap.  For each fixed old
 cap $W$, cut at its rational seams and before its extra blow-up,
 there is a constant $C_W$ such that an $S$ component satisfies
 \begin{equation}
       v_W^2\leq C_W,\qquad (\Lambda_W+v_W)^2\leq C_W.
       \label{eq:estimates:cap-squares}
 \end{equation}
 The constant depends only on $W$, its fixed line data, metric and
 boundary collars.  It is independent of the number of added
 intervals and of the new exceptional evaluations.
\item On an actually separated $N_i$ the fixed background is harmonic
 and anti-self-dual, and the metric has positive scalar curvature.
 The complementary metrics and backgrounds agree for the two bit
 choices differing by the class supported on $N_i$.
\end{enumerate}
One may first choose a positive bound for the ordinary Floer-gradient
perturbations on the finite $Y$ seams and then take those seams to any
sufficiently large finite gluing lengths.  Once those lengths have
been chosen, further compatible perturbations may be made sufficiently
small, including in total squared $L^2$ cost on the seams.  Precise
quantifiers are given in Corollary~\ref{cor:estimates:choices}.

The charge assumption concerns the sum over all retained fields and
nonnegative ideal losses.  No sign assumption is made on the charge of
an arbitrary $\PU(2)$ component.
\end{theorem}

We prove the theorem in several steps.  Only the collar, bulk-product,
and comparison-representative conclusions of
Theorem~\ref{thm:geometry:family} are used as geometric input.

\subsection{Harmonic backgrounds and local energy}

\begin{lemma}[Harmonic backgrounds and rational collars]
\label{lem:estimates:hodge}
On a completed piece with rational-homology-sphere ends, a real
cohomology class has a unique $L^2$ harmonic representative.  If pieces
are joined by increasingly long product collars with rational
cross-sections, their harmonic representatives converge on compact
subsets to the representatives of the restricted classes, provided
their total norms are bounded.  The convergence and the cutoff errors
are exponential in distance from the ends of a long collar.

Background connections can consequently be chosen to be flat on the
middle of each sufficiently long finite $Y$ collar and to obey the
product rule at all genuine faces.  On an isolated positive toric
piece their self-dual curvature is proportional to its K\"ahler form;
on a separated $N_i$ it is zero.
\end{lemma}

\begin{proof}
For clarity, ``a class'' on a completed piece means its image from
compactly supported cohomology.  The map from relative to absolute
degree-two cohomology is an isomorphism over $\R$ here: the boundary
has zero first and second real cohomology.  Minimize the two-form norm
in the affine space obtained by adding the $L^2$ closure of compactly
supported exact forms to a compactly supported representative.
Orthogonal projection gives a closed and coclosed form, which is smooth
by local elliptic regularity.  Compactly supported dual pairings
preserve its class.  For uniqueness, the difference of two such forms
is exact.  On each end it has a decreasing primitive, by the spectral
description below; a global primitive can be made equal to these
primitives outside a compact set, since $H^1$ of the cross-section
vanishes over $\R$.  Integration by parts against the difference gives
zero for its squared norm.  This is also the relative-cohomology
description of cylindrical Hodge theory
\cite[Proposition~4.9]{AtiyahPatodiSingerI}.

Here is the required quantitative collar fact.  Let $Z$ be a fixed
rational homology sphere.  Write a harmonic two-form on $[0,T]\times Z$
as $dt\wedge\alpha(t)+\gamma(t)$.  The tangential one- and two-form
Laplacians have a common positive lower spectral bound $\mu_Z$.
Each eigenmode satisfies
\[
        -f''+\mu f=0,\qquad \mu\geq\mu_Z.
\]
Its two exponential solutions and elementary integration give, for
$1\leq r\leq T/2$, and then by interior elliptic estimates,
\begin{equation}
 \|b\|_{L^2([r,T-r]\times Z)}
       \leq C_Ze^{-\sigma_Zr}\|b\|_{L^2([0,T]\times Z)},
 \qquad
 \|b\|_{C^k([r,r+1]\times Z)}
       \leq C_{Z,k}e^{-\sigma_Z\min(r,T-r-1)}\|b\|_2,
 \label{eq:estimates:rational-decay}
\end{equation}
after decreasing $\sigma_Z>0$ if necessary.  The constants can be
chosen uniformly when the link metrics range over the compact
families used in Section~\ref{sec:geometry}.  The same mode expansion
gives a primitive on an interior collar with the same bounds.  There
are no harmonic tangential modes in this argument.  In particular,
it is not an assertion about an $S^2\times S^1$ tube.

Local elliptic compactness, followed by compactly supported dual
pairings, now identifies every bounded-norm limit with the harmonic
representative of the restricted class.  To see the stronger gluing
assertion, cut off the decreasing tails of the piecewise harmonic
forms and glue their primitives on the long products.  This produces
a same-class closed representative $h_T$.  Conversely, applying the
same cutoff to the harmonic representative $b_T$ produces
representatives on the separated pieces.  Equation
\eqref{eq:estimates:rational-decay} bounds the energy change in either
direction by $Ce^{-\sigma T}$ times a bound for the total energy.
The projection identity
\begin{equation}
            \|h_T-b_T\|_2^2=\|h_T\|_2^2-\|b_T\|_2^2
            \label{eq:estimates:minimizing-remainder}
\end{equation}
therefore proves exponential closeness, with a possibly smaller
exponent.  Interior estimates give the corresponding smooth
closeness away from the splice strips.  The argument applies to
several disjoint collars at once.

For each subset of the currently long genuine collars, preglue the
harmonic data obtained by cutting just those collars.  Average these
closed representatives with the product weights formed from
$\chi(T)$ and $1-\chi(T)$, where $\chi$ is zero below one large
threshold and one above a larger threshold.  These weights depend
only on the metric parameters, so the average is closed on the
four-manifold and represents the same class.  The preceding estimate
makes its difference from harmonic data as small as prescribed.
Taking thresholds uniformly in the other disjoint collars gives the
product rule at every multiple face.  At a $J$ face the data on a
piece depend only on that piece's parameters.  At an $N_i$ face the
piece is already separated, so its datum remains exactly harmonic.
The thresholds may be increased as the auxiliary sizes vary; thus
the errors can have any prescribed smaller order in the estimates
below.

A line bundle on a rational cross-section admits the flat connection
specified by its torsion class; no trivialization of that line bundle
is asserted.  The small primitive in
\eqref{eq:estimates:rational-decay} cuts a connection off to this flat
one.  Connections with the prescribed curvatures exist because their
periods are the prescribed integral periods.  Fixed collar gauges and
the fixed line identifications give the asserted compatibility.

Finally, in four dimensions the star on two-forms and their $L^2$ norm
are conformally invariant.  The toric K\"ahler form has finite squared
norm before completion, since the moment polygon has finite area.
It therefore remains an $L^2$ self-dual harmonic form after conformal
cylindrical completion.  The main piece has $b^+=1$, so this form
spans its self-dual harmonic space.  Since $b^+(N_i)=0$, its harmonic
datum is anti-self-dual.  A class supported on $N_i$ restricts to zero
on the complement; the product construction accordingly gives the
same exterior datum for the two bit choices.
\end{proof}

We next record the local estimate with the error norms that will
actually be used.  The equations and Clifford convention are those of
Section~\ref{sec:indices}; equivalently,
\[
 D_a\Phi=P_D,\qquad
 \rho(F_a^{0,+})=(\Phi\Phi^*)_{00}+P_+.
\]
These are the usual trace-free monopole equations
\cite{FeehanLenessPU2I}.  If $\Phi$ is identified with a two-by-two
matrix with singular values $s,t$, then
\[
 \bigl\langle(\Phi\Phi^*)_{00}\Phi,\Phi\bigr\rangle
 =\frac14(s^4+t^4)+\frac52s^2t^2
 \geq\frac14|\Phi|^4.
\]
We write $d_*=1/4$ for this lower bound; changing the common quadratic
normalization merely changes this fixed positive constant.

\begin{lemma}[Local integrated estimate]
\label{lem:estimates:local}
Set
\[
 V=\frac{\Scal}{4}\id+\frac12\rho(b),\qquad
 h=\max\{0,-\lambda_{\min}(V)\}.
\]
If $\zeta=f^2$ is a compactly supported real cutoff with $0\leq f\leq1$,
then
\begin{equation}
 \begin{split}
 \frac12\int\zeta^2|\nabla_a\Phi|^2
   +\frac{d_*}{4}\int\zeta^2|\Phi|^4
 \leq{}&2\int\zeta^2|P_D|^2
       +\frac1{d_*}\int\zeta^2(h^2+|P_+|^2)\\
       &+\frac9{d_*}\int\frac{|d\zeta|^4}{\zeta^2}.
 \end{split}
 \label{eq:estimates:cutoff}
\end{equation}
The last integrand is defined continuously as $16|df|^4$ where
$f=0$.  Enlarging the numerical constants in
\eqref{eq:estimates:cutoff}, if the curvature perturbation is expressed
as a two-form instead of an endomorphism, accounts for its fixed
Clifford norm.
In particular, uniform local geometry, bounded $h$, and bounded
local squared error costs give uniform bounds for
$\int(|\nabla_a\Phi|^2+|\Phi|^4)$ on interior unit strips.  No derivative
of $P_D$ occurs.
\end{lemma}

\begin{proof}
Pair the weak identity $D_a^*D_a\Phi=D_a^*P_D$ with
$\zeta^2\Phi$, and use the positive-spinor Weitzenb\"ock formula.
The right side is
\[
 \langle P_D,D_a(\zeta^2\Phi)\rangle
 =\zeta^2|P_D|^2+
      2\zeta\langle P_D,\rho(d\zeta)\Phi\rangle.
\]
Thus no integration by parts differentiates the prescribed error.
The covariant derivative cross term is at most
$\frac12\zeta^2|\nabla_a\Phi|^2+2|d\zeta|^2|\Phi|^2$;
the displayed Dirac cross term is at most
$\zeta^2|P_D|^2+|d\zeta|^2|\Phi|^2$.
The negative zeroth-order terms are bounded by
$\zeta^2(h+|P_+|)|\Phi|^2$.  Apply Young's inequality separately to
these two terms, using a total of at most $d_*/2$ of the quartic
coefficient, and to $3|d\zeta|^2|\Phi|^2$, using at most $d_*/4$.
The resulting remainders are bounded by the right side of
\eqref{eq:estimates:cutoff}.  Finally
$|d(f^2)|^4/(f^2)^2=16|df|^4$.  A cutoff supported in a fixed larger
strip and equal to one on the strip in question proves the last
assertion.
\end{proof}

\begin{lemma}[Charge on a finite $Y$ middle]
\label{lem:estimates:cylinder}
Consider a flat-background product cylinder with fixed cross-section
$Y$.  On its translation-invariant middle, assume that the ordinary
connection perturbation is the gradient of a functional
$\mathcal W$ with
$\|\mathcal W\|_\infty\leq M$ and uniformly bounded gradient.
Let $r,e$ denote any additional connection and Dirac errors there.
Suppose their total squared $L^2$ norms are bounded by $\nu$.
There is a constant $C_Y$, depending only on the fixed cross-section
and the functional and gradient bounds, independent of both $\nu$
and cylinder length, such that one can
choose boundary slices in prescribed interior unit strips near its
two ends for which
\begin{equation}
                  \kappa_{\rm middle}\geq-C_Y-C_Y\nu.
                  \label{eq:estimates:cylinder-charge}
\end{equation}
The complementary bounded strips may be included in the retained
region of a cutoff estimate.
\end{lemma}

\begin{proof}
Use temporal gauge and the convention
$\rho_3(\xi)=-\rho_4(dt\wedge\xi)$ on positive spinors.  Write the
equations as
\[
 \dot\Phi+D_a\Phi=e,\qquad
 \dot a+*F_a=q(\Phi)+\sigma\nabla\mathcal W(a)+r,
 \qquad \sigma\in\{1,-1\}.
\]
With the stated Clifford and trace conventions $c_*=1/2$, and
$\langle q(\Phi),v\rangle=-c_*\langle D'_a(v)\Phi,\Phi\rangle$.
All pairings in this proof are integrated over $Y$ and are real
pairings.  For $H(t)=\langle D_a\Phi,\Phi\rangle$, differentiation
and the Dirac equation give the exact identity
\begin{equation}
 \begin{split}
 -\langle *F_a,\dot a\rangle
 ={}&\|\dot a\|_2^2+c_*H'(t)+2c_*\|\dot\Phi\|_2^2
       -\sigma\frac{d}{dt}\mathcal W(a)\\
    &-\langle r,\dot a\rangle-2c_*\langle e,\dot\Phi\rangle.
 \end{split}
 \label{eq:estimates:cylinder-identity}
\end{equation}
In particular an ordinary Floer perturbation contributes endpoint
values, rather than a length times its pointwise size.  Young's
inequality absorbs the last two terms into half the positive energies
at a cost of $\frac12\|r\|_2^2+c_*\|e\|_2^2$.

It remains to bound $H$ on suitable slices.  Apply
Lemma~\ref{lem:estimates:local} on fixed larger strips around the two
prescribed unit strips.  Their geometry and background are fixed, and
the gradient has a uniform pointwise bound, so
\[
 \int_{J\times Y}(|\nabla_a\Phi|^2+|\Phi|^4)\leq A_Y(1+\nu).
\]
Since $|D_a\Phi|\leq\sqrt3|\nabla_a\Phi|$ for the tangential Dirac
operator,
\[
 \int_J|H(t)|\,dt
 \leq\sqrt3\,\Vol(J\times Y)^{1/4}
          \{A_Y(1+\nu)\}^{3/4}.
\]
There is a slice in each unit strip with
$|H|\leq B_Y(1+\nu)^{3/4}$, where $B_Y$ is independent of $\nu$.
Integrating
\eqref{eq:estimates:cylinder-identity} between the chosen slices and
using
\[
 \kappa_{[u,v]\times Y}
       =-\frac1{4\pi^2}\int_u^v\langle *F_a,\dot a\rangle\,dt
\]
gives the explicit lower bound
\[
 -\frac{2c_*B_Y(1+\nu)^{3/4}+2M+
             \frac12\|r\|_2^2+c_*\|e\|_2^2}{4\pi^2}.
\]
Since $(1+\nu)^{3/4}\leq1+\nu$, this proves
\eqref{eq:estimates:cylinder-charge} for every $\nu\geq0$ with the
asserted constant.  The same local
estimate controls the adjoining cutoff strips.  Notice that the
proof does not require pointwise differentiability bounds for $e$.
\end{proof}

\subsection{The long tube and the small-width test}

\begin{proposition}[Tube comparison and retained-region energy]
\label{prop:estimates:energy}
Fix $\mathcal P$.  Let $\epsilon\downarrow0$ through the family in
Section~\ref{sec:geometry}, allowing all admissible deeper
regularizations and rational isolations.  Let $\mathcal R_\epsilon$
be the union of the retained regions, obtained by removing the long
translation-invariant $Y$ middles.  Keep initial $Y$ prefixes of
length $R_0=C\epsilon^{-1/2}$, as well as the bounded cutoff strips.
For harmonic backgrounds, and also for the product-rule backgrounds
of Lemma~\ref{lem:estimates:hodge} chosen with smaller errors,
\begin{equation}
 \sum_{\mathcal T}\|\beta-\lambda_{\mathcal T}\omega_{S^2}\|_{L^2(\mathcal T)}^2
   +\|\beta\|_{L^2(\mathrm{rest})}^2=o_{\mathcal P}(\epsilon^{-1}),
 \qquad \|\beta\|_2^2=O_{\mathcal P}(\epsilon^{-1}).
 \label{eq:estimates:tube-background}
\end{equation}
Here $\mathcal T$ runs over the bulk tubes, bounded end pieces may be
transferred to the rest, and $\omega_{S^2}$ is the area form of
$S^2(1/\sqrt2)$.  The subscript on the little-$o$ means uniformity
over those admissible choices and all parameter faces for the fixed
problem.

For solutions whose total charge is bounded above, bounded ordinary
$Y$ gradient perturbations and additional errors of squared cost
$o(\epsilon^{-1})$ give
\begin{equation}
 \int_{\mathcal R_\epsilon}
       (|\Phi|^4+|F_a^{0,+}|^2+|F_a^0|^2)
                  =o_{\mathcal P}(\epsilon^{-1}).
 \label{eq:estimates:small-energy}
\end{equation}
For each fixed $\epsilon>0$, the same quantities have a finite bound
$E_{\mathcal P,\epsilon}$ independent of all subsequently lengthened
rational collars and of the omitted $Y$ lengths, with the error
costs bounded.  The assertions also hold for the retained main fields
of ideal tuples with nonnegative omitted losses.
\end{proposition}

\begin{proof}
The geometric input is Lemma~\ref{lem:geometry:comparison}.  After
collapsing long products to bounded products, prescribe the vertical
form $\lambda_{\mathcal T}\omega_{S^2}$ on their middle slices and zero
real curvature on rational middle slices.  The restriction classes
agree with these prescriptions.  Relative de Rham representatives on
the bounded truncations, followed by reinsertion of the products,
give a same-class comparison form $h$.  Uniform bounded truncations
at the smooth pieces, conical annuli and rescaled satellites give
\begin{equation}
 \|h\|_2^2
  =\sum_{\mathcal T}\lambda_{\mathcal T}^2\Vol(\mathcal T)
       +o_{\mathcal P}(\epsilon^{-1}).
 \label{eq:estimates:comparison-upper}
\end{equation}
In this formula the volumes are those of the product models.  Metric
errors change its two sides by $o_{\mathcal P}(\epsilon^{-1})$.
Additional rational cylinders have zero comparison curvature on
their long parts, so their arbitrary lengths add no leading term.

For the converse estimate, on each exact product fiber
$\Sigma=S^2(1/\sqrt2)$ one has
\[
       \int_\Sigma\beta=2\pi\lambda_{\mathcal T},\qquad
       \operatorname{area}(\Sigma)=2\pi.
\]
Fiberwise Cauchy--Schwarz and integration over the other two
coordinates therefore bound the energy from below by the summand
in \eqref{eq:estimates:comparison-upper}.  More precisely, the cross
term against the model vertical form is fixed by this same flux, so
on the exact product
\[
 \|\beta-\lambda_{\mathcal T}\omega_{S^2}\|_2^2
       =\|\beta\|_2^2-
                    \lambda_{\mathcal T}^2\Vol(\mathcal T).
\]
First discard bounded tip ranges and an arbitrarily small end
fraction of each long tube.  On the remainder the metric and its
needed derivatives tend uniformly to the product model.  The errors
in the last identity are bounded by this metric error times the
energy, which is $O_{\mathcal P}(\epsilon^{-1})$ by minimization and
\eqref{eq:estimates:comparison-upper}.  Then let the discarded
fraction tend to zero.  Nonnegativity of all the remaining energies
proves \eqref{eq:estimates:tube-background}.

The background is flattened on each $Y$ middle only after the
prefix of length $R_0$.  Equation
\eqref{eq:estimates:rational-decay} bounds the added squared error
by $C_{\mathcal P}e^{-2\sigma R_0}\epsilon^{-1}$, which is negligible.
The product-rule alterations at genuine faces can have smaller
cost by the threshold choice in Lemma~\ref{lem:estimates:hodge}.

On a bulk tube, $\Scal=4$ in the model and
\[
 \frac{\Scal}{4}=1\geq\frac{|\lambda_{\mathcal T}|}{2}
   =\frac12\|\rho(i\lambda_{\mathcal T}\omega_{S^2})\|,
 \qquad |\lambda_{\mathcal T}|\leq2.
\]
Consequently the squared negative part $h^2$ from
Lemma~\ref{lem:estimates:local} has integral
$o_{\mathcal P}(\epsilon^{-1})$ there, by
\eqref{eq:estimates:tube-background} and the integrated metric-error
bounds.  The same assertion holds on the rest: its background
energy is lower order, its negative scalar-curvature squared cost
is lower order by the geometric decomposition, and the retained
$Y$ prefixes have bounded local scalar cost times
$O(R_0)=o(\epsilon^{-1})$.  A bounded $Y$ gradient perturbation has
the same lower-order squared cost on those prefixes.  All other
errors have the stipulated total cost.  Thus, for a cutoff equal
to one on the retained part away from the bounded transition strips,
\[
 \int_{\mathcal R_\epsilon}(h^2+|P_+|^2+|P_D|^2)
                        =o_{\mathcal P}(\epsilon^{-1}).
\]
Apply \eqref{eq:estimates:cutoff}; use the local strip estimate at its
finitely many transitions.  At a genuine infinite end use cutoffs
and finite-energy tails.  This gives the $|\Phi|^4$ assertion.
The quadratic curvature equation gives the assertion for
$|F_a^{0,+}|^2$.

For the full curvature, sum charges over all retained components.
If $Q$ is the fixed upper charge bound and there are $q_Y$ omitted
$Y$ middles, Lemma~\ref{lem:estimates:cylinder} gives
\begin{equation}
       \kappa(\mathcal R_\epsilon)
       \leq Q+q_YC_Y(1+\nu).
       \label{eq:estimates:retained-charge}
\end{equation}
Nonnegative particles and other omitted nonnegative losses only
decrease the left side.  Substitution into
\eqref{eq:estimates:energy-charge} proves the full-curvature
assertion.  This sums over the arbitrary $\PU(2)$ fields as well as
the reduced fields, and never discards a component on an assumed
charge sign.

Finally fix $\epsilon$.  All bounded-truncation contributions in the
same argument are then finite uniform constants; lengthening a
positive rational cylinder adds neither negative scalar cost nor
a zero-mode background.  The omitted-middle bound is independent
of length.  Repeating the estimates with constants instead of
little-$o$ proves the stated $E_{\mathcal P,\epsilon}$ bound.
\end{proof}

\begin{corollary}[The small-width test]
\label{cor:estimates:tube}
For $\epsilon$ sufficiently small, every $R$ or $S$ reduction on a
small-width transition or replacement region has $vU=0$.
This choice is uniform under the subsequent deeper regularizations
and rational isolations allowed by the metric construction.
\end{corollary}

\begin{proof}
Such a region contains a bulk tube of length at least
$c_{\mathcal P}/\epsilon$, whose other circle factor has length
bounded below.  For a reduction, $F_a^0$ has diagonal entries
$iD/2,-iD/2$, so $|F_a^0|^2=|D|^2/2$.  Each fiber has
$\int_{S^2}D=2\pi u$.  Cauchy--Schwarz therefore gives, with a
constant independent of the auxiliary lengths,
\[
        \int_{\mathcal T}|F_a^0|^2
                       \geq c'_{\mathcal P}u^2/\epsilon.
\]
If the integer $u$ is nonzero this contradicts
\eqref{eq:estimates:small-energy} along a sequence
$\epsilon\downarrow0$.  The estimates are uniform in precisely the
other parameters permitted in such a sequence.  The stationary
overlap of the small-width and toric prescriptions is included.
\end{proof}

\subsection{Period inequalities and finite persistence}

\begin{lemma}[The period inequality on an isolated positive piece]
\label{lem:estimates:projection}
Let $M$ be a positive-width toric main limit, with period class
$H=U+aX_1+b_0X_2$, $a,b_0\in[0,1/4]$, $a+b_0>0$.
Let the fixed background be harmonic.  Every finite-energy solution
of the unperturbed reduced equations satisfies
\begin{equation}
        vH=0\quad\hbox{or}\quad
        (vH)((\Lambda+v)H)<0,
        \label{eq:estimates:period}
\end{equation}
with $vH=0$ when the spinor vanishes and the strict alternative when
it is nonzero.
\end{lemma}

\begin{proof}
At a reduction the trace-free curvature is
$\diag(iD/2,-iD/2)$.  The two bundle-diagonal blocks of the doubly
trace-free quadratic term are likewise one half of the active
spinor quadratic term with opposite signs.  The reduced curvature
equation is therefore
\[
             \rho(iD^+)=(\psi\psi^*)_0,
\]
and the active Dirac operator has fixed-plus-variable curvature
$i(\beta+D)$.  This checks that $v$ and $\Lambda+v$ in
\eqref{eq:estimates:period} are full line classes.
For the Clifford convention in use,
\[
 \tr\bigl(\rho(i\eta^+)\rho(i\xi^+)\bigr)
       =4\langle\eta^+,\xi^+\rangle,
 \qquad |\psi|^4=8|D^+|^2.
\]
The integrated Dirac identity consequently reads
\begin{equation}
 0=\int_M\left(|\nabla\psi|^2+\frac{\Scal}{4}|\psi|^2\right)
       +2\int_M\left(\langle\beta^+,D^+\rangle+|D^+|^2\right).
 \label{eq:estimates:line-Weitzenbock}
\end{equation}
To justify integration on the completion, insert cutoffs changing
on bounded-width end strips.  The derivative error there is
bounded by a fixed constant times
$\Vol(\mathrm{strip})^{1/2}
 (\int_{\mathrm{strip}}|\psi|^4)^{1/2}$ and tends to zero.
The other terms follow by the finite curvature energy and the
cutoff estimate.  In particular no global $L^2$ assumption on
$\psi$ is needed.

Scalar curvature is nonnegative and positive on a nonempty interior
open set.  The first integral in
\eqref{eq:estimates:line-Weitzenbock} is therefore nonnegative.
If it vanished, then $\nabla\psi=0$ and
$\Scal|\psi|^2=0$ everywhere.  The latter equality makes $\psi$
vanish on that open set, and a parallel section which vanishes at
a point vanishes on the connected manifold $M$.  Thus the first
integral is strictly positive whenever $\psi\ne0$.  Write
$D^+=D_H+D_\perp$, where $D_H$ is its self-dual harmonic projection.
By Lemma~\ref{lem:estimates:hodge}, $\beta^+$ belongs to the same
one-dimensional harmonic space.  With $h_H$ the harmonic form of
class $H$,
\[
 \beta^+=2\pi\frac{\Lambda H}{H^2}h_H,
 \qquad D_H=2\pi\frac{vH}{H^2}h_H.
\]
The second integral in \eqref{eq:estimates:line-Weitzenbock} is
therefore
\[
 (2\pi)^2\frac{(vH)((\Lambda+v)H)}{H^2}
                         +\|D_\perp\|_2^2.
\]
It is negative if $\psi\ne0$, which proves the strict alternative.
If $\psi=0$, the curvature equation gives $D^+=0$ and hence $vH=0$.
\end{proof}

\begin{lemma}[Vertex test and finite persistence]
\label{lem:estimates:vertices}
After $\epsilon$ has been fixed as in
Corollary~\ref{cor:estimates:tube}, sufficiently small regularization,
sufficiently long finite rational isolation, and sufficiently small
perturbations give \eqref{eq:estimates:clamp} on every positive bridge.
The assertion is uniform near genuine faces and over all sufficiently
long finite $Y$ seams.
\end{lemma}

\begin{proof}
\emph{The isolated vertex test.}
First consider an isolated positive-width limit.  Put
$p=4a$, $q=4b_0$.  Its period is the convex combination
\begin{equation}
 H=(1-p)(1-q)H_\varnothing+p(1-q)H_1
              +(1-p)qH_2+pqH_{12}.
 \label{eq:estimates:convex}
\end{equation}
All weights are nonnegative, and at least one nonempty vertex has
positive weight.  All nonempty vertices have negative
$\Lambda$-evaluation by \eqref{eq:estimates:vertex-values}; the empty
one has $\lambda\leq0$.  Consequently $\Lambda H<0$.
If $u\ne0$, its parity gives
\begin{equation}
                         |u|\geq|\lambda|.
                         \label{eq:estimates:parity-size}
\end{equation}
Indeed this is immediate for $\lambda=0,-1$, and for $\lambda=-2$
a nonzero even integer has absolute value at least two.

For type $R$, if all available nonempty vertices violated the
asserted weak inequality, multiplication of
\eqref{eq:estimates:convex} by $\operatorname{sign}(u)v$ would give
$\operatorname{sign}(u)vH>0$, contradicting
Lemma~\ref{lem:estimates:projection}.  For type $S$, failure at all
such vertices, together with \eqref{eq:estimates:parity-size}, gives
\[
       \operatorname{sign}(u)vH
       \geq\sum_I\theta_I|\lambda_I|
       =|\Lambda H|>0.
\]
Here $\theta_I$ are the four weights.  A number of magnitude at
least $|\Lambda H|$ cannot satisfy either alternative of
\eqref{eq:estimates:period}.  This proves the limiting test.

\emph{Persistence at finite sizes.}
Fix $\mathcal P$ and $\epsilon$, and suppose the conclusion failed
along a sequence
with regularization tending to zero, isolations tending to infinity,
and perturbations tending to zero.  Pass to a fixed sign of $u$ and
to limiting period parameters.  The small-width alternative is already
excluded by Corollary~\ref{cor:estimates:tube}.
By Proposition~\ref{prop:estimates:energy} the retained energy is
uniformly bounded.  We will pass to an unperturbed reduction on the
retained smooth main piece.  Only the pairings of vertices with positive
limiting weight must pass to this limit; those vertices remain supported
on that piece even when other classes collapse.

On a ball, Abelian Coulomb gauge gives a
connection one-form bounded in $W^{1,2}$ by its curvature norm.
Lemma~\ref{lem:estimates:local}, followed by
$d\psi=\nabla_a\psi-a\psi$ and $W^{1,2}\subset L^4$, bounds the
spinor in $W^{1,2}$ as well.  To exclude an atom of its fourth-power
measure at a point, take a radial cutoff $f_R$ equal to one on
$B_{R^2}$ and zero outside $B_R$, with logarithmic transition.
In four dimensions
\[
          \int|df_R|^4\leq C|\log R|^{-3}.
\]
Use $\zeta=f_R^2$ in \eqref{eq:estimates:cutoff}, first pass to the
sequence limit and then let $R\downarrow0$.  The bounded background
cost is $O(R^4)$ and the error costs tend to zero, so the atom is
zero.  The critical Sobolev concentration argument now gives strong
local $L^4$ convergence of the spinors: its defect measure is
supported at atoms, as follows by applying the Sobolev inequality
to a cutoff times the weakly convergent difference and partitioning
sets on which the limiting gradient-energy measure is small.
The reduced curvature equation and the interior elliptic estimate
for $d^*\oplus d^+$ then give strong $W^{1,2}$ convergence of the
Coulomb connection forms on smaller balls.

For completeness, bootstrap on balls on which their $L^4$ norms
are small.  In the Dirac equation multiplication by this one-form
is a small operator from $W^{1,p}$ to $L^p$, for $2<p<4$, and can
be absorbed in the local Dirac estimate.  The cutoff inhomogeneity
is already in $L^p$; choosing $p>8/3$ gives spinor integrability
greater than eight.  The quadratic curvature equation then gives
connection coefficients in $W^{1,q}$ for some $q>4$.  Repeated
Dirac and $d^*\oplus d^+$ estimates give smooth convergence when
the errors tend to zero in the corresponding local Sobolev norms.
This proves convergence to an unperturbed reduction without
differentiating the errors in the initial energy estimate or using
transversality.  Ideal markings elsewhere and filled true lens ends
do not change its compact pairings.

Only vertices of positive limiting
weight in \eqref{eq:estimates:convex} matter.  By
Lemma~\ref{lem:geometry:period-support}, each such vertex is
represented rationally in the smooth retained main part: it omits
a side whose low coefficient has vanished, and includes a collapsed
high divisor with coefficient $1/4$, making it orthogonal to that
divisor.  Compactly supported dual representatives therefore bound
its evaluation by the retained $L^2$ energy.  Those evaluations lie
in a fixed discrete lattice, so a subsequence makes them constant.
All of their nonempty sides are whole for late terms of the
sequence.  There is no need to bound an evaluation on a vertex of
zero limiting weight, nor to bound $u$ if the empty weight is zero.

For $R$, strict violation on a nonempty vertex persists as a
strictly positive discrete value; the empty vertex, if relevant,
has $\operatorname{sign}(u)u>0$.  For $S$, violation is the closed
inequality $\operatorname{sign}(u)z_I\geq|\lambda_I|$, and persists
as well.  The limiting calculation just made is a contradiction,
including if an $S$ spinor limits to zero.  The argument allowed
arbitrary genuine faces and arbitrary further $Y$ lengths, proving
the stated uniformity.
\end{proof}

\subsection{Fixed caps and the order of choices}

\begin{proposition}[Cap constants before exterior choices]
\label{prop:estimates:cap}
For each old cap $W$ there is a metric, unchanged on its rational
end collars, and a number $C_W<\infty$ such that
\[
 \begin{gathered}
 \forall\,\mathcal P\ \exists\,L_{\mathcal P},\delta_{\mathcal P}>0:\\
 \text{isolation at least }L_{\mathcal P},\quad
 \text{admissible errors at most }\delta_{\mathcal P}\\
 \Longrightarrow\quad
 \text{no cap-containing }R\text{ and \eqref{eq:estimates:cap-squares}.}
 \end{gathered}
\]
Here $\mathcal P$ is any fixed finite exterior problem whose
restriction to $W$ is the fixed cap datum.  Neither $C_W$ nor that
metric depends on the new exceptional evaluation or the number of
intervals.  The thresholds and intermediate full-curvature bounds
may depend on both.
\end{proposition}

\begin{proof}
The cap contains a nonzero class $A$ with $A^2=0$ and $cA$ odd.
Its real intersection form is nondegenerate, since its boundaries
are rational homology spheres.  A nondegenerate negative-definite
form has no nonzero isotropic vector, so $b^+(W)>0$.
Moreover every allowed restriction $v_W\equiv c_W\pmod2$ has odd
evaluation on $A$, and is nonzero in real cohomology.

Choose the metric so that none of these countably many integral
classes is anti-self-dual.  Here is the period argument with the
permitted support.  Suppose a nonzero harmonic representative
$\alpha$ of one such class is anti-self-dual.  Choose a nonzero
self-dual harmonic form $\omega$.  Unique continuation for harmonic
forms~\cite[Remark~3]{Aronszajn1957} implies
that the nonzero sets of $\alpha$ and $\omega$ meet in any prescribed
open region for metric variation.  Infinitesimal changes of the
metric there give all homomorphisms between the anti-self-dual and
self-dual two-form spaces.  Thus one may choose a compactly
supported variation of the star for which
\[
                       \int_W\langle\dot *\alpha,\omega\rangle\ne0.
\]
Differentiating the harmonic projection shows that this is, up to
a nonzero fixed factor, its derivative in the $\omega$ direction.
The wall for this class is therefore closed with empty interior
among metrics fixed on the end collars.  Baire's theorem avoids
all the walls at once.  This is the usual local period variation
argument; see also \cite{DonaldsonKronheimer1990}.  It uses the
nonzero real class, so torsion presents no exception.

Fix this metric and its harmonic datum $\beta_W$.  Let
$\omega_1,\ldots,\omega_b$ be an $L^2$ orthonormal basis of its
self-dual harmonic forms.  They and their end primitives decrease
exponentially by \eqref{eq:estimates:rational-decay}.  Cover $W$ by
finitely many interior sets and bounded-overlap unit end strips
$V_k$, $k\geq0$.  On a completed-cap limit, the local estimate gives
\begin{equation}
                 \|D^+\|_{L^2(V_k)}\leq A_W
                 \quad\hbox{for every }k,
                 \label{eq:estimates:cap-local}
\end{equation}
where $A_W$ depends only on the fixed cap geometry and background.
Indeed the scalar curvature and background are bounded on all
these sets, the enlarged sets have uniformly bounded volume and
cutoffs, and the limiting equations are unperturbed.  No incoming
spinor boundary value occurs in Lemma~\ref{lem:estimates:local}.
For $R$ the left side is zero.  Hence the convergent sums
\begin{equation}
 B_{W,a}=\frac{A_W}{2\pi}\sum_k\|\omega_a\|_{L^2(V_k)},
 \qquad
 B_W=\left(\sum_{a=1}^bB_{W,a}^2\right)^{1/2}
 \label{eq:estimates:cap-constant}
\end{equation}
bound the coefficients, and then the norm, of the limiting positive
projection of $v_W$.  These numbers contain no exterior data.

We must identify that projection with the topological restriction,
without imposing a cap-only bound on the anti-self-dual curvature.
Fix $\mathcal P$ and $\epsilon$ first.  Compactly supported dual
pairings in $W$ and Proposition~\ref{prop:estimates:energy} bound all
integral evaluations of $v_W$; thus only finitely many restriction
classes occur for this problem.  In a hypothetical sequence with
isolation tending to infinity, pass to a constant one of these
classes.  Local convergence gives a field on the completed cap.

For a harmonic $\omega_a$, let $\alpha_a$ be its exponentially
decreasing primitive on each end.  If $\chi_N$ is one through
distance $N$ and zero beyond $N+1$, then
\[
             \widehat\omega_{a,N}
                =\chi_N\omega_a+d\chi_N\wedge\alpha_a
\]
is a compactly supported closed representative of its relative
class.  The sign follows from $d\alpha_a=\omega_a$; equivalently it
is $\omega_a-d((1-\chi_N)\alpha_a)$ on the ends.
Its pairing with $D/(2\pi)$ is the topological pairing of $v_W$
with $[\omega_a]$.

The full-curvature bound needed in these end strips can depend on
$\mathcal P$.  Retain a prefix of length $N+2$ of a neighboring $Y$
collar and choose its omitted-middle slice beyond the prefix.
Lemma~\ref{lem:estimates:cylinder} still gives the same lower
middle-charge bound.  Lemma~\ref{lem:estimates:local} bounds the
self-dual cost of each additional cap-prefix strip by a fixed
constant.  Equation~\eqref{eq:estimates:energy-charge} then gives
\begin{equation}
 \int_{\text{cap and prefixes through }N+2}|D|^2
            \leq E_{\mathcal P,\epsilon}+C'_W(N+2).
 \label{eq:estimates:cap-prefix}
\end{equation}
The same bound applies at a positive rational isolation collar.
For each fixed $N$ it follows either before passing to the cap
limit, once that prefix is isolated, or afterward.  Its intercept
may be arbitrarily large as exterior data vary.
The correction pairing is consequently bounded by
\begin{equation}
 C_W' e^{-\sigma N}
              \sqrt{E_{\mathcal P,\epsilon}+C'_W(N+2)}
                         \longrightarrow0
 \quad(N\longrightarrow\infty),
 \label{eq:estimates:cap-tail}
\end{equation}
with $\mathcal P$ fixed.  The remaining integral converges by
\eqref{eq:estimates:cap-local} and exponential decay of $\omega_a$.
It equals $(2\pi)^{-1}\int_W\langle D^+,\omega_a\rangle$.
This proves the required identification and the bound
$\|v_W^+\|\leq B_W$.

An $R$ limit would have $v_W^+=0$, contrary to the chosen period
condition.  For an $S$ limit, possibly with vanishing limiting
spinor,
\[
 v_W^2\leq\|v_W^+\|^2\leq B_W^2,\qquad
 (\Lambda_W+v_W)^2
       \leq(\|\Lambda_W^+\|+B_W)^2.
\]
Choose once and for all
\begin{equation}
 C_W=1+\max\{B_W^2,(\|\Lambda_W^+\|+B_W)^2\}.
 \label{eq:estimates:CW}
\end{equation}
For the fixed problem, a sequence violating either finite square
bound has a constant restriction class and would contradict this
strict cushion in the limit.  A sequence of cap-containing $R$
fields similarly contradicts the period condition.  Thus finite
thresholds exist.  The exceptional negative summand is outside
$W$ under the rational orthogonal splitting, and has not been
included in either cap constant.
\end{proof}

\begin{corollary}[Order of finite choices]
\label{cor:estimates:choices}
The constants $C_W$ can be fixed before choosing the delays,
exceptional evaluations and repetition count in the stack.  After
those choices have specified a finite problem $\mathcal P$, there
are $\epsilon>0$, admissible finite auxiliary sizes, numbers
$\delta>0$ and $T_0<\infty$, with the following property.  For every
collection of finite $Y$ lengths $T_j\geq T_0$ and ordinary
Floer-gradient perturbations of size at most $\delta$, there is a
positive allowed size for the further non-gradient perturbations
such that all conclusions of Theorem~\ref{thm:estimates:clamp} hold.
The further size may depend on the chosen finite $T_j$; it includes
their total squared error costs.  At genuine ends the perturbations
and backgrounds obey the product rule and the specified decreasing
tail bounds.
\end{corollary}

\begin{proof}
First fix the cap metrics and \eqref{eq:estimates:CW}.  Next fix
$\mathcal P$, including any arbitrarily large exceptional
evaluations.  Choose $\epsilon$ by
Corollary~\ref{cor:estimates:tube}.  Bounded $Y$ gradients only cost
$O(\epsilon^{-1/2})$ on its retained prefixes and contribute bounded
endpoint terms on the omitted middles, so this first choice remains
valid under the later lengthenings.  Choose the smaller
regularization parameter $d$, the rational isolation lengths and
the perturbation bounds for Lemma~\ref{lem:estimates:vertices} and
Proposition~\ref{prop:estimates:cap}, with this $\epsilon$ fixed.

Here is the uniformity argument for the order of the last choices.
If no such finite auxiliary choices, positive $\delta$, and $T_0$
worked for all $T_j\geq T_0$, choose failures successively with
$d\to0$, each required isolation length tending to infinity,
gradient sizes tending to zero, and $\min T_j\to\infty$.
For each chosen finite collection of lengths, let the additional
errors tend to zero in the indicated local and total norms before
selecting a failure.  The off-middle energy bound is uniform in
all these lengths.  After passing to subsequences, the small-width
test, the positive-width vertex argument or the cap argument
therefore supplies exactly the contradiction already proved.
This is a neighborhood assertion in all the auxiliary variables:
it supplies a fixed positive gradient tolerance before the finite
$Y$ gluing lengths are selected.  It is stronger than merely
finding one successful diagonal sequence.

Finally fix those finite lengths.  Further errors can be decreased
in norms depending on all the sizes now chosen.  Compactly
supported errors have finite cost; errors on genuine collars can
have prescribed square-integrable decreasing tails.  For an error
on a finite $Y$ middle this permits, for example, a pointwise bound
whose square times the chosen length tends to zero.  Alternatively
that middle may remain exactly a gradient equation.  The argument
never requires a fixed non-gradient pointwise tolerance on
arbitrarily long uncontrolled middles.  On a separated positive
$N_i$ allow only the small compact instanton perturbations that
preserve its spinor-zero conclusion.  This completes the choice
of all data and the proof of Theorem~\ref{thm:estimates:clamp}.
\end{proof}
